\begin{definition}[Normal one-site separation hypotheses]%
    \label{def:peps_normalOneSiteSeparationHypotheses}
    \lean{TNLean.PEPS.NormalOneSiteSeparationHypotheses}
    \leanok
    \uses{def:peps_regionInjectivityData, def:peps_regionComplement}
    For each site, the normal theorem assumes two injective regions with
    injective complements that differ exactly at that site. This is the
    one-site comparison used after the edge gauges have been obtained.
    Diagrammatically, in the square-lattice proof this is the comparison
    between the source regions $R$ and $S=R\cup\{v\}$:
    \begin{tenkzequation}
        % Source verdict: paper_normal.tex:1519--1544 compares the injective
        % regions R and S=R\cup\{v\}; their complements are injective as well.
        % Type verdict: this is a lattice/set schematic, not a tensor
        % contraction. The marked site v remains present with all its bonds.
        % Cell-set audit: R=(2-4,2-4)-(4,4) has eight selected cells, while
        % S=(2-4,2-4) has nine; v=(4,4) is marked in both pictures.
        % Contraction/boundary audit: there are no tensor contractions and no
        % tensor boundary signature; the boundary legs are cut lattice edges.
        \tenkzkernel
        \begin{tenkz}[rows={ket,ket,ket,ket,ket}, cols=5,
            west=open, east=open, north=open, south=open]
            \tn[at=(4,4), species=marked, name=vSite, label pos=sw]{v}
            \tnmark[form=enclosure, slot=selected, tint, label pos=nw]
                {(2,2) .. (4,4) - (4,4)}{$R$}
        \end{tenkz}
        $\qquad$
        \begin{tenkz}[rows={ket,ket,ket,ket,ket}, cols=5,
            west=open, east=open, north=open, south=open]
            \tn[at=(4,4), species=marked, name=vSite, label pos=sw]{v}
            \tnmark[form=enclosure, slot=selected, tint, label pos=nw]
                {(2,2) .. (4,4)}{$S$}
        \end{tenkz}
        $\qquad R\subset S=R\cup\{v\}$
    \end{tenkzequation}
\end{definition}

\begin{theorem}[Membership in a one-site separated region]%
    \label{thm:peps_normalOneSiteSeparation_mem_withSite_iff}
    \lean{TNLean.PEPS.NormalOneSiteSeparationHypotheses.mem_withSite_iff}
    \leanok
    \uses{def:peps_normalOneSiteSeparationHypotheses}
    For the two regions associated to a site $v$, a vertex $w$ belongs to
    the region containing $v$ if and only if either $w=v$, or $w$
    belongs to the comparison region omitting $v$.
\end{theorem}
\begin{proof}\leanok
    If $w=v$, use the distinguished-site membership and nonmembership
    assumptions. If $w\neq v$, use the equality of the two regions away
    from the distinguished site.
\end{proof}

\begin{definition}[Normal square-lattice blocking regions]%
    \label{def:peps_normalSquareBlockingRegions}
    \lean{TNLean.PEPS.NormalSquareBlockingRegions}
    \lean{TNLean.PEPS.normalSquareBlockingRegions_of_TCover}
    \leanok
    \uses{def:peps_normalRectangleInjectivityHypotheses,
        def:peps_regionInjectivityData,
        lem:peps_normalSquareRegionRS_injective_of_unionClosure,
        lem:peps_normalSquareRegionT_injective_of_cover}
    In the translationally invariant square-lattice proof, the source paper
    assumes that every $2\times 3$ and $3\times 2$ region is injective.
    The formal statement here is the corresponding abstract, scope-restricted
    hypothesis bundle: it records the rectangular injectivity assumptions, the
    size bound $n,m\geq 7$ with $|V|=nm$, and three finite regions
    $R$, $S$, and $T$ whose injectivity is assumed directly. It does
    not by itself assert that these regions are the displayed square-lattice
    unions, nor that their translated variants give the edge-centred blocking
    around every edge. Conversely, for the coordinate square lattice with
    $n,m\geq7$, the rectangular hypotheses, union closure, and a rectangular
    cover of the displayed $T$-region yield this abstract $R,S,T$
    structure, with $R$ and $S$ obtained from their rectangular
    decompositions and $T$ obtained from the supplied cover.
    % Current source-comparison status: docs/paper-gaps/peps_normal_ft_section3_route.tex.
    The source regions are:
    \begin{tenkzequation}
        % Source/formula verdict: Molnar et al., arXiv:1804.04964,
        % paper_normal.tex:1405--1429 and source PDF p. 15 depict
        % R=(2-4,2-4)\{(4,4)} and S=(2-4,2-4).
        % Ink-to-index verdict: every dot is an ordinary tensor; every lattice
        % edge is a virtual index, and the selected fill records membership only.
        % Structural count verdict (not mechanically emitted as a lattice
        % signature): each 5-by-5 panel has 25 sites, 40 internal virtual
        % contractions, and 20 cut virtual bonds; (4,4) remains ordinary.
        % Structurally derived boundary/type verdict: no physical index is
        % exposed, so each panel has (V,P)=(20,0).
        \tenkzkernel
        \begin{tenkz}[rows={ket,ket,ket,ket,ket}, cols=5,
            west=open, east=open, north=open, south=open]
            \tnmark[form=enclosure, slot=selected, tint, label pos=nw]
                {(2,2) .. (4,4) - (4,4)}{$R$}
        \end{tenkz}
        $\qquad\text{and}\qquad$
        \begin{tenkz}[rows={ket,ket,ket,ket,ket}, cols=5,
            west=open, east=open, north=open, south=open]
            \tnmark[form=enclosure, slot=selected, tint, label pos=nw]
                {(2,2) .. (4,4)}{$S$}
        \end{tenkz}
    \end{tenkzequation}
    \begin{tenkzequation}
        % Formula verdict: in row-column coordinates the exact local formal
        % region is T=(1-6,1-5) minus the vertical block (3-5,1-2) and
        % the shifted horizontal block (2-3,3-5).
        % Ink-to-index verdict: every dot is an ordinary tensor, every lattice
        % edge is a virtual contraction, and the selected fill records only
        % membership in T; it adds no contraction or tensor operation.
        % Contraction/boundary verdict: the 6-by-5 panel has 49 internal
        % virtual contractions and 22 cut virtual bonds, with no exposed
        % physical indices, hence boundary type (V,P)=(22,0).
        % Source verdict: arXiv:1804.04964, paper_normal.tex:1430--1445
        % depicts the same 5-by-6 window complement; this panel uses the exact
        % cell set of normalSquareRegionT rather than the source's rounded ink.
        \tenkzkernel
        \begin{tenkz}[rows={ket,ket,ket,ket,ket,ket}, cols=5,
            west=open, east=open, north=open, south=open]
            \tnmark[form=enclosure, slot=selected, tint, label pos=n]
                {(1,1) .. (6,5) - (3,1) .. (5,2) - (2,3) .. (3,5)}{$T$}
        \end{tenkz}
    \end{tenkzequation}
\end{definition}

\begin{theorem}[Normal edge blocking to a three-site injective chain]%
    \label{thm:peps_normalEdgeBlockingToInjectiveChain}
    \notready
    \uses{lem:peps_injectiveRegion_union,
        def:peps_normalSquareBlockingRegions,
        def:peps_normalEdgeBlockingHypotheses,
        thm:peps_normalEdgeBlockingHypotheses_injectiveChain,
        thm:peps_edgeBlockedInsertionAlgebraIsomorphism}
    Under the square-lattice normality hypotheses of
    \cite[Theorem~3]{Molnar2018NormalPEPS}, the regions $R$, $S$, and
    $T$, together with their translated variants, give an edge-centred
    blocking into red, blue, and complementary injective regions. These are
    the regions on which the three-site injective-chain insertion argument is
    applied:
    \begin{tenkzequation}
        % Formula verdict: A1=(4-6,2-3), A2=(3-4,4-6), and A3 is their
        % complement in the 7-by-7 frame; blocking contracts each region to
        % one injective chain site and carries the marked A1--A2 edge across.
        % Ink-to-index verdict: left-hand dots are ordinary tensors and lattice
        % edges are virtual contractions; right-hand boxes are the three
        % blocked tensors, joined by virtual bonds with physical legs exposed.
        % Contraction/boundary verdict: the lattice has 84 internal virtual
        % contractions and boundary type (V,P)=(28,0); the cyclic blocked
        % network has no open virtual bonds and three physical legs, so
        % (V,P)=(0,3).
        % Source verdict: exact discrete translation and contraction summary of
        % arXiv:1804.04964, paper_normal.tex:1475--1499.
        \tenkzkernel
        \begin{tenkz}[rows={ket,ket,ket,ket,ket,ket,ket}, cols=7,
            west=open, east=open, north=open, south=open]
            % The authored marked wire restates the (4,3)--(4,4) grid bond.
            \tnwire[species=marked]{(4,3)}{(4,4)}
            \tnmark[form=enclosure, slot=selected, tint]
                {(4,2) .. (6,3)}{$A_1$}
            \tnmark[form=enclosure, slot=secondary, tint]
                {(3,4) .. (4,6)}{$A_2$}
            \tnmark[form=enclosure, slot=complement, label pos=n]
                {(1,1) .. (7,7) - (4,2) .. (6,3) - (3,4) .. (4,6)}{$A_3$}
        \end{tenkz}
        $\quad\Rightarrow\quad$
        \begin{tenkz}[rows={ket}, cols=3, physical=up, boundary=periodic]
            \tn[skin=box]{A_1} & \tn[skin=box]{A_2} & \tn[skin=box]{A_3}
            \tnwire[species=marked]{(1,1)}{(1,2)}
            \tnmark[form=enclosure, slot=selected]{(1,1)}{}
            \tnmark[form=enclosure, slot=secondary]{(1,2)}{}
            \tnmark[form=enclosure, slot=complement]{(1,3)}{}
        \end{tenkz}
    \end{tenkzequation}
\end{theorem}

\begin{theorem}[Translationally invariant normal gauge absorption]%
    \label{thm:peps_tiNormalGaugeAbsorption}
    \notready
    \uses{thm:peps_normalEdgeBlockingToInjectiveChain,
        thm:peps_edge_gauge_absorption}
    In the translationally invariant normal square-lattice case, the edge
    gauges obtained after blocking are represented by one horizontal matrix
    $X$ and one vertical matrix $Y$, unique up to scalar. Absorbing these
    matrices into $B$ gives a tensor $\widetilde B$, equivalently the final
    local gauge formula of \cite[Theorem~3]{Molnar2018NormalPEPS}:
    % Formula (paper_normal.tex:1451--1471):
    % B = lambda (Y on north)(Y^{-1} on south)(X on west)(X^{-1} on east) A.
    % Ink-to-index: the four gauge capsules act on the four virtual indices of A;
    % the diagonal north-east stub is the unchanged physical index.
    % Boundary signature: both sides have (V,P)=(4,1).  Source verdict: exact
    % translation-invariant square-lattice gauge absorption in house style.
    \begin{tenkzequation}
        \tenkzkernel
        % Each unbonded typed port renders its own stub, so the four virtual
        % faces and the diagonal physical index need no boundary atoms.
        \begin{tenkz}[rows={wire}, cols=1, bonds=none]
            \tn[label pos=sw,
                ports={180:virtual, 0:virtual, 90:virtual, 270:virtual,
                       45:physical}]{B}
        \end{tenkz}
        $ = \lambda $
        \begin{tenkz}[rows={wire}, cols=1, bonds=none]
            \tn[name=tiAbsA, label pos=sw,
                ports={180:virtual, 0:virtual, 90:virtual, 270:virtual,
                       45:physical}]{A}
            \tn[at=1 w of tiAbsA, name=tiAbsX, skin=ring, species=operator,
                ports={0:virtual, 180:virtual}]{X}
            \tn[at=1 e of tiAbsA, name=tiAbsXinv, skin=ring,
                species=operator,
                ports={180:virtual, 0:virtual}]{X^{-1}}
            \tn[at=1 n of tiAbsA, name=tiAbsY, skin=ring, species=operator,
                ports={270:virtual, 90:virtual}]{Y}
            \tn[at=1 s of tiAbsA, name=tiAbsYinv, skin=ring,
                species=operator,
                ports={90:virtual, 270:virtual}]{Y^{-1}}
            \tnwire{tiAbsA.180}{tiAbsX.0}
            \tnwire{tiAbsA.0}{tiAbsXinv.180}
            \tnwire{tiAbsA.90}{tiAbsY.270}
            \tnwire{tiAbsA.270}{tiAbsYinv.90}
        \end{tenkz}
    \end{tenkzequation}
    The torus form of this absorption, with the gauge family carried onto
    every edge by the translations from one matrix per orientation class, is
    part of Theorem~\ref{thm:fundamentalTheorem_normalTorusPEPS_unconditional}.
    The present open-lattice entry awaits a translation-invariance predicate
    on the open lattice, which would make the per-edge gauges of the edge
    blocking one horizontal and one vertical matrix.
\end{theorem}

\begin{theorem}[Fundamental Theorem for normal square-lattice PEPS]%
    \label{thm:fundamentalTheorem_normalSquarePEPS}
    \notready
    \uses{thm:peps_tiNormalGaugeAbsorption,
        thm:peps_oneVertexComplementComparison,
        def:GaugeEquivPEPS}
    Let $A$ and $B$ be translationally invariant normal square-lattice PEPS
    tensors such that every $2\times 3$ and $3\times 2$ region is
    injective. If they generate the same state on an $n\times m$ region with
    $n,m\geq 7$, then they are related by horizontal and vertical gauges
    $X,Y$, up to a scalar $\lambda$ satisfying
    $\lambda^{nm}=1$. This is
    \cite[Theorem~3]{Molnar2018NormalPEPS}.
    Theorem~\ref{thm:fundamentalTheorem_normalSquarePEPS_unconditional} proves
    an interior-window form of this statement from the source hypotheses alone,
    without translation invariance and with per-vertex scalars in place of the
    single $\lambda$; the gauge-equivalence conclusion on a general connected
    graph, from region injectivity alone, is
    Theorem~\ref{thm:fundamentalTheorem_normalPEPS_connected}. The present
    source-exact entry awaits a translation-invariance predicate on the open
    lattice, which would make the per-edge gauges one horizontal and one
    vertical matrix and the per-vertex scalars a single $\lambda$ with
    $\lambda^{nm}=1$.
\end{theorem}

\begin{theorem}[Fundamental Theorem for normal square-lattice PEPS,
        interior-window form]%
    \label{thm:fundamentalTheorem_normalSquarePEPS_unconditional}
    \lean{TNLean.PEPS.fundamentalTheorem_normalSquarePEPS_unconditional}
    \leanok
    \uses{thm:fundamentalTheorem_normalSquarePEPS,
        def:peps_normalRectangleInjectivityHypotheses,
        def:applyGauge,
        def:peps_edgeInsertedCoeff}
    Let $A$ and $B$ be PEPS tensors on the open $n\times m$ square lattice
    with the same positive physical dimension, equal and positive bond
    dimensions, generating the same state, and such that every contiguous
    $2\times 3$ and $3\times 2$ coordinate rectangle of either tensor is
    injective. Then there is a family of invertible matrices, one on each
    edge, such that:
    \begin{enumerate}
        \item at every edge with interior margins (for a horizontal edge with
              left endpoint $(x,y)$: $3\le x$, $x+7\le n$, $4\le y$,
              $y+6\le m$; for a vertical edge the mirrored margins), inserting
              any matrix on the edge of $A$ gives the same coefficient as
              inserting it on the corresponding edge of the gauge-absorbed $B$,
              for every physical configuration; and
        \item every vertex $v=(x,y)$ of the interior comparison window
              $4\le x\le n-9$, $4\le y\le m-9$ carries a nonzero scalar
              $\lambda_v$ such that $A_v$ equals $\lambda_v$ times the gauge
              action of the family on $B$ at $v$, on every virtual configuration
              of the incident edges and every physical index.
    \end{enumerate}
    This is the open-lattice form of \cite[Theorem~3]{Molnar2018NormalPEPS},
    restricted in two ways. First, the conclusion reaches only the interior:
    the absorbed equality holds at the edges with interior margins and the
    per-vertex relation at the window vertices, while vertices near the
    lattice boundary are not reached, the blocking frames of the open lattice
    requiring interior margins. Second, the scalars of distinct vertices are
    not asserted equal, and no condition $\lambda^{nm}=1$ is asserted: the
    single $\lambda$ of the source comes from translation invariance, which is
    not assumed here, and without it a single scalar is false --- rescaling
    the components of $A$ at two interior vertices by $\mu$ and $\mu^{-1}$
    preserves the state and the rectangular injectivity but separates the two
    scalars. No lower bound on $n$ or $m$ is assumed; both conclusions are
    empty when the lattice is too small to contain interior edges or window
    vertices.
    \begin{proof}
        \leanok
        \uses{lem:peps_injectiveRegion_union_pos}
        The union closure of injective regions follows from the positive bond
        dimensions and Lemma~\ref{lem:peps_injectiveRegion_union_pos}. The
        translated blocking frames at the edges with interior margins produce
        per-edge gauges, whose absorption into $B$ gives the absorbed
        inserted-coefficient equality at every such edge. A gauge preserves
        blocked-region linear independence, so at every window vertex the
        region comparison applies to the gauge-absorbed tensor: comparing the
        $3\times 3$ square at the vertex minus its corner with its one-site
        completion yields two proportionality scalars $c_R$ and $c_S$, every
        boundary edge of either region carrying interior margins, and the
        inserted-site extraction gives the per-vertex relation with the ratio
        $\lambda_v = c_S/c_R$.
    \end{proof}
\end{theorem}
